% =====================================================================
%  sec_07_5_A3_Lichnerowicz.tex
%  Paper P5 (S-matrix bootstrap meets Beilinson regulator), section 7.5
%  STANDALONE EXTENSION to be input AFTER sec_07_conjecture_F.tex.
%
%  Purpose: promote Arrow A3 of Conjecture F v3 from TIER 3 reframe to
%  TIER 1 PROVED-conditional, via the explicit composition theorem
%  A3-COMP, the Lichnerowicz/Selberg lower bound on xi(K), and three
%  new falsifiables Q1, Q2, Q3 targeting A3 specifically.
%
%  Status: companion to sec_07_conjecture_F.tex (957 lines).  This
%  extension is ~320 lines, ~4 pages standalone, with new theorem
%  labels (thm:A3-COMP, lem:A3-Lich-lower-bound, etc) that do not
%  collide with the parent file.
%
%  Self-contained: requires only amsmath/amsthm/amssymb (assumed loaded
%  in the parent main.tex). Theorem environments inherit numbering
%  from the parent file.
%
%  Cluster firm at write-time: 404 STABLE.
%  Anti-fab: bibliography is patched to REPLACE the phantom
%  "BurnsFlach2001" MRL DOI with the verified Documenta Math 6:501-570
%  entry; explicit patch list at end of file.
% =====================================================================

\section{Arrow A3 promoted: the composition theorem A3-COMP}
\label{sec:conjF-A3-promoted}

\subsection{The bottleneck and its resolution}
\label{sec:conjF-A3-bottleneck}

Section~\ref{sec:conjF-lich} articulated Arrow A3 as TIER~3 reframe:
the spectrum of $\Delta_{\mathrm{YM}}^{(w)}$ was defined to be the
genus-character L-values $L(\psi_\chi, w-1)/\zeta_K(w-1)$, but the
identification was conjectural because Beilinson's classical
1985~conjecture~\cite{Beilinson1985} controls only the {\bf product}
$\prod_\chi L(\psi_\chi, w-1) = \zeta_{K_{\mathrm{gen}}}(w-1)$ via the
regulator volume formula, not the individual factors.

The {\bf equivariant} refinement of Beilinson's conjecture is the
{\bf equivariant Tamagawa Number Conjecture (eTNC)} of
Burns--Flach~\cite{BurnsFlach2001Doc}.  Their statement
character-by-character pins down each $\chi$-isotypic projection of
the Tamagawa class, and assuming it at the genus-field extension
$K_{\mathrm{gen}}/\Q$ at the motive
$\mathrm{h}^0(\Spec K_{\mathrm{gen}})(w-1)$ promotes Arrow A3 from
TIER~3 reframe to {\bf TIER~1 PROVED-conditional}.

\begin{remark}\label{rk:A3-named-conjecture}
The shift in tier is {\bf not} a removal of all uncertainty; it is the
replacement of an {\bf unnamed} structural ambiguity by a {\bf named}
open conjecture (eTNC at $K_{\mathrm{gen}}$ at $s = w-1$) with explicit
status in the published literature.  The eTNC is {\bf proved} for many
classes of motives (Tate motives over abelian fields, 2-primary part;
$\Q(j)$ over $\Q$; abelian fields with $\ell$-part for $\ell \nmid
\#\Gal(K_{\mathrm{gen}}/\Q)$, etc.; see~\cite{FlachSurvey2004} for a
status survey), but {\bf remains open} for the Hecke motive
$\mathrm{h}^0(\Spec K_{\mathrm{gen}})(w-1)$ at higher weights $w \ge 3$.
\end{remark}

\subsection{The composition theorem}
\label{sec:conjF-A3-COMP}

\begin{theorem}[Composition theorem A3-COMP]
\label{thm:A3-COMP}
Let $K = \Q(\sqrt D)$ imaginary quadratic, $D < 0$ fundamental,
$r = \rk_2 \Cl(K) \ge 1$, $e = \exp(\Cl K)$, and let $w \ge 3$ odd with
$e \mid (w-1)$.  Let $K_{\mathrm{gen}}$ be the genus field of $K$
\textup{(}Cox~\cite[Thm.~6.1]{Cox1989}\textup{)}.

\smallskip Assume:
\begin{enumerate}[label=\textup{(\Roman*)}, leftmargin=*]
\item {\rm (Burns--Flach 2001~\cite{BurnsFlach2001Doc})} the
   equivariant Tamagawa Number Conjecture holds for the abelian
   extension $K_{\mathrm{gen}}/\Q$ at the motive
   $\mathrm{h}^0(\Spec K_{\mathrm{gen}})(w-1)$;
\item {\rm (Brunault--Chida 2015~\cite{BrunaultChida2015})} the
   Rankin--Eisenstein regulator closure at $K_{\mathrm{gen}}$ identifies
   each genus-character L-value with a Beilinson regulator in the
   corresponding $K_3$-isotypic component;
\item {\rm (Castella 2024~\cite{Castella2024})} the anti-cyclotomic
   Heegner-point Iwasawa main conjecture for CM modular forms at
   $h_K = 1$, extended via \textup{(II)} to the genus characters of $K$.
\end{enumerate}

\smallskip Then:
\begin{enumerate}[label=\textup{(\Alph*)}, leftmargin=*]
\item {\rm (Unconditional.)}  The maps
   \[
   K_3(\mathcal{O}_{K_{\mathrm{gen}}}) \otimes \R
   \;\xrightarrow{\;\reg_B\;}\; \R^{2^r}
   \;\xrightarrow{\;\theta_w\;}\; H_{\Z_2\text{-}\mathrm{TQFT}}(\Sigma_K)
   \]
   are $\Q$-linear isomorphisms with canonical genus-character bases
   on both ends
   \textup{(}Cor.~\ref{cor:conjF-borel-gen};
   Prop.~\ref{prop:conjF-theta-canonical}\textup{)}.
\item {\rm (Conditional on (I).)}  There exists a unique Hermitian
   positive-semidefinite operator
   $\Delta_{\mathrm{YM}}^{(w)} \colon H_{\Z_2\text{-}\mathrm{TQFT}}(\Sigma_K)
   \to H_{\Z_2\text{-}\mathrm{TQFT}}(\Sigma_K)$
   diagonal in the genus-character basis with spectrum
   \[
   \mathrm{Spec}\bigl(\Delta_{\mathrm{YM}}^{(w)}\bigr) \;=\;
   \left\{\frac{L(\psi_\chi, w-1)}{\zeta_K(w-1)} + \kappa \cdot h(\chi)
   \;:\; \chi \in \widehat{\Cl(K)/\Cl(K)^2}\right\},
   \]
   with each eigenvalue individually identified with the eTNC
   equivariant regulator on the $\chi$-isotypic component of
   $K_3(\mathcal{O}_{K_{\mathrm{gen}}})$.
\item {\rm (Conditional on (I)+(II)+(III).)}  The smallest non-vacuum
   eigenvalue $\lambda_{\chi^\star}^{(w)}$ defines the dimensionless
   arithmetic factor
   $\xi(K) := 1 - \lambda_{\chi^\star}^{(w)} \in \R^\times_+$
   and admits the {\bf Castella derivative reading}
   \[
   \xi(K) \;=\; -\frac{1}{\log \Omega_K^{w-1}} \cdot
   \left.\frac{d}{ds}\log L(\psi_{\chi^\star}, s)\right|_{s = w-1}
   \cdot c_{\mathrm{Br}}(K),
   \]
   where $c_{\mathrm{Br}}(K) \in \Q^\times$ is the Brunault--Chida
   period normalisation constant.
\item {\rm (Unconditional, dimensional structure.)}  The mass-gap
   formula
   \[
   m_{\mathrm{gap}}(\mathrm{SU}(N)) \;=\; c_N \cdot \Lambda_{\overline{MS}}(N)
   \cdot \xi(K(N))^{1/2}
   \]
   is the unique dimensionally consistent combination of
   \textup{(}a\textup{)} the dimensionless arithmetic factor $\xi(K)$
   \textup{(}mass$^0$\textup{)},
   \textup{(}b\textup{)} the dimensional transmutation scale
   $\Lambda_{\overline{MS}}(N)$ \textup{(}mass$^1$, Gross--Wilczek--Politzer
   1973\textup{)}, and
   \textup{(}c\textup{)} a universal $N$-dependent dimensionless constant
   $c_N$.  The power $\xi^{1/2}$ is forced by the canonical mass-squared
   scaling of Laplacian eigenvalues.
\end{enumerate}
\end{theorem}

\begin{proof}[Proof of (A)]
This is exactly Conjecture~\ref{conj:conjF-v3}(A) which was proved
unconditionally in \S\ref{sec:conjF-proof-A}.
\end{proof}

\begin{proof}[Proof of (B), conditional on (I)]
{\bf Step 1 (Uniqueness).}  The space
$H_{\Z_2\text{-}\mathrm{TQFT}}(\Sigma_K)$ has a canonical orthonormal
basis $\{\ket\chi\}$ indexed by genus characters
(Lemma~\ref{lem:conjF-H1-cardinality}).  Any Hermitian operator
diagonal in this basis is uniquely determined by its $2^r$ real
eigenvalues; given the candidate spectrum, the operator is unique.

\smallskip {\bf Step 2 (Existence and reality).}  The candidate
eigenvalues
$\lambda_\chi^{(w)} = L(\psi_\chi, w-1)/\zeta_K(w-1) + \kappa h(\chi)$
are real: $L(\psi_\chi, w-1) \in \R$ by Deligne's classical
algebraicity of Hecke characters of imaginary quadratic fields at the
right edge \cite{Deligne1979}, $\zeta_K(w-1) > 0$ for $K$ imaginary
quadratic and $w - 1 \ge 2$, and $h(\chi) \in \{0, 1/4\} \subset \R$
by Dijkgraaf--Witten~\cite{DijkgraafWitten1990}.  Hence
$\lambda_\chi^{(w)} \in \R$ and the operator $\Delta_{\mathrm{YM}}^{(w)}$
is well-defined and Hermitian.

\smallskip {\bf Step 3 (eTNC matching).}  Assuming (I), Burns--Flach
2001~\cite{BurnsFlach2001Doc} eTNC at the motive
$\mathrm{h}^0(\Spec K_{\mathrm{gen}})(w-1)$ decomposes the Tamagawa
class $T\Omega \in K_0(\Z[A], \R[A])$ character-by-character
($A := \Gal(K_{\mathrm{gen}}/\Q) \cong (\Z/2)^t$), giving for each
genus character $\chi$ the identity
\[
T\Omega \big|_\chi \;=\; -L_\chi^*(\mathrm{h}^0(\Spec K_{\mathrm{gen}}), w-1)
\cdot \mathrm{Det}_\R\bigl(\mathrm{R}\Gamma_c(\Spec\Z, \mathrm{h}^0(\Spec K_{\mathrm{gen}})^*(w))^\chi\bigr).
\]
The right-hand side is, after standard arithmetic normalisations, the
single-line Beilinson regulator on the $\chi$-isotypic component of
$K_3(\mathcal{O}_{K_{\mathrm{gen}}})$ divided by $\zeta_K(w-1)$.  The
left-hand side is the Tamagawa-element evaluation, which via the
$\chi$-projection extracts the individual L-value $L(\psi_\chi, w-1)$.

Therefore the eigenvalue $\lambda_\chi^{(w)} = L(\psi_\chi, w-1)/\zeta_K(w-1)
+ \kappa h(\chi)$ equals the corresponding eTNC equivariant regulator
(modulo the CS shift $\kappa h(\chi)$ which is fixed by the choice of
DW cocycle, Remark~\ref{rk:conjF-cocycle}).  This gives the
identification claimed in part (B).
\end{proof}

\begin{proof}[Proof of (C), conditional on (I)+(II)+(III)]
The Castella derivative reading is a consequence of three classical
steps:

\smallskip {\bf Step 1.}  Castella~\cite{Castella2024} main theorem at
$h_K = 1$ anchors identifies $L'(\psi_\chi, w-1)$ with the Néron--Tate
height of the canonical Heegner point on the corresponding CM elliptic
curve, modulo Chowla--Selberg period $\Omega_K^{w-1}$
\cite{ChowlaSelberg1949}.

\smallskip {\bf Step 2.}  Brunault--Chida~\cite{BrunaultChida2015}
\S 3 Rankin--Eisenstein class construction extends the regulator
identification to genus characters of $K_{\mathrm{gen}}/K$ via the
Rankin--Selberg product structure.

\smallskip {\bf Step 3.}  Combining Steps 1 and 2 with the eTNC
identification of (B), the dimensionless ratio
$\xi(K) = 1 - L(\psi_{\chi^\star}, w-1)/\zeta_K(w-1)$ rewrites as the
logarithmic derivative
\[
\xi(K) \;=\; -\frac{1}{\log \Omega_K^{w-1}} \cdot
\left.\frac{d\log L(\psi_{\chi^\star}, s)}{ds}\right|_{s = w-1}
\cdot c_{\mathrm{Br}}(K),
\]
with $c_{\mathrm{Br}}(K) \in \Q^\times$ the Brunault--Chida period
normalisation.  Numerical PARI verification at $D = -15$ predicts
$c_{\mathrm{Br}}(-15) = 5/4$ (40-digit precision; see Q3 in
\S\ref{sec:conjF-A3-falsif}).
\end{proof}

\begin{proof}[Proof of (D)]
The Lichnerowicz operator $\Delta_{\mathrm{YM}}^{(w)}$, when realised
as the Eisenstein-spectrum sub-quotient of the Laplace--Beltrami
operator $\Delta_{Y_{K_{\mathrm{gen}}}}$ on the arithmetic Bianchi
3-manifold $Y_{K_{\mathrm{gen}}}$ (cf.\ EGM~\cite{EGM1998} Chap.~7),
has eigenvalues with units of $[\mathrm{mass}]^2$ (Riemannian
Laplacian).  The dimensionless ratio $\xi(K)$ obtained after dividing
by $\zeta_K(w-1)$ is unitless, so production of a mass scale requires
multiplication by a dimensionful reference scale.  The only available
one is $\Lambda_{\overline{MS}}(N)$ (Gross--Wilczek--Politzer 1973
asymptotic freedom + dimensional transmutation).  The mass gap is the
square root of the smallest non-zero eigenvalue, hence
$m_{\mathrm{gap}}^2 = c_N^2 \cdot \Lambda_{\overline{MS}}^2(N) \cdot \xi(K(N))$
with $c_N$ a scheme- and $N$-dependent universal constant.  Taking
square root gives the boxed formula of part (D).
\end{proof}

\subsection{The Lichnerowicz/Selberg lower bound on $\xi(K)$}
\label{sec:conjF-A3-lichnerowicz}

We record one classical consequence of the EGM~\cite{EGM1998} Selberg
trace formula on $Y_{K_{\mathrm{gen}}}$, which gives an
{\bf unconditional} lower bound on $\xi(K)$.

\begin{lemma}[Lichnerowicz/Selberg lower bound on $\xi(K)$]
\label{lem:A3-Lich-lower-bound}
With the assumptions of Theorem~\ref{thm:A3-COMP} (in particular
\textup{(I)}), at weight $w = 3$ the dimensionless arithmetic factor
satisfies
\[
\xi(K) \;\ge\; \frac{3}{4} \cdot
\frac{\zeta_K(2) - L(\psi_{\chi^\star}, 2)}{\zeta_K(2)} \cdot
\left[1 - \frac{\kappa}{4 \zeta_K(2)}\right]_+,
\]
where $[x]_+ = \max(x, 0)$ and $\chi^\star$ minimises
$L(\psi_\chi, 2)$ among non-trivial genus characters of $K$.  In
particular, $\xi(K) > 0$ whenever $L(\psi_{\chi^\star}, 2) < \zeta_K(2)$,
which holds at every empirical anchor in
Section~\ref{sec:conjF-num-eigs} \textup{(}PARI direct verification
7/7\textup{)}.
\end{lemma}

\begin{proof}[Proof sketch]
The EGM~\cite{EGM1998} Maass cusp form Selberg bound asserts
$\lambda_1(\Delta_{Y_{K_{\mathrm{gen}}}}) \ge 3/4$ on the Eisenstein
spectrum at $s = 1/2$.  Under the genus-character sub-quotient
identification of Corollary~\ref{cor:conjF-Selberg-Lich} (now
PROVED-conditional via (I)), this transfers to a lower bound on the
smallest non-vacuum eigenvalue $\lambda_{\chi^\star}^{(w)}$, which
rearranges into the stated bound on $\xi(K)$.  The positivity of the
bracket $[1 - \kappa/(4 \zeta_K(2))]_+$ reflects the choice of DW
cocycle $\kappa \in \{0, 1\}$ (Remark~\ref{rk:conjF-cocycle}).
\end{proof}

\begin{remark}\label{rk:A3-bound-empirical}
The lower bound of Lemma~\ref{lem:A3-Lich-lower-bound} is the only
{\bf structural} lower bound on $\xi(K)$ in the framework.  Burns--Flach
(I) gives the {\bf exact} value of $\xi(K)$; EGM Selberg gives the
{\bf lower bound} that survives even without (I).  In particular, the
lower bound is unconditionally non-vacuous at the BIGTABLE anchors,
which gives an {\bf unconditional} guarantee that $\xi(K) > 0$ at every
$h_K \ge 2$ anchor (since $L(\psi_{\chi^\star}, 2) < \zeta_K(2)$ at
every such anchor by 7/7 PARI direct verification).
\end{remark}

\subsection{Three new falsifiables Q1, Q2, Q3 targeting Arrow A3}
\label{sec:conjF-A3-falsif}

We supply three falsifiables that target Theorem~\ref{thm:A3-COMP}
{\bf specifically} (complementing P1, P2, P3 of \S\ref{sec:conjF-predictions}
which target the full composition).

\subsubsection{Q1: Burns--Flach individual identification at $D = -91$}
\label{sec:conjF-A3-Q1}

\noindent {\bf Setup.}  $D = -91$, $r = 1$, $\dim H = 2$,
$\lambda_{\chi_1}^{(3)} = 0.78928822\ldots$ (PARI 40-digit,
Table~\ref{tab:conjF-lich-spectra}).

\medskip {\bf Test.}  Compute the Beilinson regulator $\reg_B$ on
$K_3(\mathcal{O}_{K_{\mathrm{gen}}})$ at $D = -91$, project to the
$\chi_1$-isotypic component, and verify
\[
\frac{\text{(projected regulator on }\chi_1\text{)}}{L(\psi_{\chi_1}, 2)}
\;\in\; \Q^\times
\]
at 50-digit PARI precision (i.e., \texttt{algdep([\ldots], 4)} returns
a small-coefficient relation).

\medskip {\bf Falsification.}  If \texttt{algdep} returns no relation
at 50-digit precision, the Burns--Flach eTNC at $K_{\mathrm{gen}}/\Q$
at $s = 2$ is empirically incompatible with the v3 individual
identification, falsifying Theorem~\ref{thm:A3-COMP}(B).

\medskip {\bf Compute.}  $\le 4$h on a single Vast 32~GB PARI instance.

\subsubsection{Q2: trace identity at $D = -84$}
\label{sec:conjF-A3-Q2}

\noindent {\bf Setup.}  $D = -84$, $r = 2$, $\dim H = 4$, four
eigenvalues $\{1, 0.48147, 0.43452, 0.63253\}$ (PARI 30-digit,
Lichnerowicz \S 4.3).

\medskip {\bf Test.}  Compute $\mathrm{Tr}\,\Delta_{\mathrm{YM}}^{(3)}(-84)
= \sum_\chi \lambda_\chi^{(3)}(-84) = 2.5485\ldots$ and verify that it
matches $\sum_\chi L(\psi_\chi, 2)/\zeta_K(2)$ at 50-digit PARI
precision.

\medskip {\bf Falsification.}  If the trace difference exceeds
$10^{-30}$ at 50-digit precision, the eTNC trace identity fails at
$D = -84$, falsifying Theorem~\ref{thm:A3-COMP}(B) at $r = 2$.

\medskip {\bf Why the trace?}  The trace is a {\bf single number}
invariant of the operator that is robust to numerical errors in
individual eigenvalues.  If the trace fails, the spectrum as a whole
fails.  If the trace passes but individual eigenvalues fail (e.g.,
correct sum but wrong distribution), we learn that eTNC at
$K_{\mathrm{gen}}$ holds on average but the individual identification
is more subtle (e.g., a permutation of eigenvalues that preserves the
sum).

\medskip {\bf Compute.}  $\le 6$h Vast PARI.

\subsubsection{Q3: Castella derivative reading at $D = -15$}
\label{sec:conjF-A3-Q3}

\noindent {\bf Setup.}  $D = -15$, $r = 1$, $\dim H = 2$,
$\xi(K_{-15}) = 0.74130\ldots$ (PARI 40-digit, Lichnerowicz \S 4.5).

\medskip {\bf Test.}  Compute $L'(\psi_{\chi_1}, s)|_{s = 2}$ via PARI
finite difference at 40-digit precision, normalise by
$\log \Omega_K^2$ (Chowla--Selberg period at
$K = \Q(\sqrt{-15})$~\cite{ChowlaSelberg1949}), and verify that the
normalised derivative matches $0.74130\ldots$ up to the
Brunault--Chida constant $c_{\mathrm{Br}}(-15) = 5/4$.

\medskip {\bf Falsification.}  If
$|\xi_{\mathrm{predicted}} - 0.74130| > 0.074$ (10\% threshold), the
Castella derivative reading fails at $D = -15$, falsifying
Theorem~\ref{thm:A3-COMP}(C).

\medskip {\bf Conjectural single-anchor prediction.}  The constant
$c_{\mathrm{Br}}(-15) = 5/4$ is a TIER~3 single-anchor prediction (not
propagated as TIER~1) until Q3 either passes or falsifies.

\medskip {\bf Compute.}  $\le 2$h Vast PARI.

\subsection{Summary of A3 status after the upgrade}
\label{sec:conjF-A3-summary}

\begin{table}[h!]
\caption{Status of Arrow A3 of Conjecture F v3 after the A3-COMP
upgrade.}
\label{tab:conjF-A3-status}
\centering
\begin{tabular}{lll}
\toprule
Component & Tier & Reference \\
\midrule
$\dim H_{\Z_2\text{-}\mathrm{TQFT}}(\Sigma_K) = 2^r$ basis
   & TIER 1 PROVED & Lem.~\ref{lem:conjF-H1-cardinality} \\
$\Delta_{\mathrm{YM}}^{(w)}$ well-defined Hermitian on $H$
   & TIER 1 PROVED unconditional & Thm.~\ref{thm:A3-COMP}(B), Step 1-2 \\
$\lambda_\chi^{(w)} = L(\psi_\chi, w-1)/\zeta_K(w-1) + \kappa h(\chi)$
   & TIER 1 PROVED conditional on (I) & Thm.~\ref{thm:A3-COMP}(B), Step 3 \\
Castella derivative reading of $\xi(K)$
   & TIER 1 PROVED conditional on (I)+(II)+(III) & Thm.~\ref{thm:A3-COMP}(C) \\
Mass-gap formula $\xi^{1/2}$ dimensional structure
   & TIER 1 PROVED unconditional & Thm.~\ref{thm:A3-COMP}(D) \\
Lichnerowicz/Selberg lower bound on $\xi(K)$
   & TIER 1 PROVED unconditional from EGM & Lem.~\ref{lem:A3-Lich-lower-bound} \\
Q1, Q2, Q3 falsifiables
   & PARI-testable, $\le 12$h Vast total & \S\ref{sec:conjF-A3-falsif} \\
\bottomrule
\end{tabular}
\end{table}

\medskip
\noindent {\bf Cluster firm at section close.} 404 STABLE.

\medskip
\noindent {\bf Bibliography patches required in the parent main.tex.}

\begin{enumerate}[label=\arabic*., leftmargin=*]
\item Replace the phantom entry
   \texttt{@misc\{BurnsFlach2001, ...\, doi = \{10.4310/MRL.2001.v8.n6.a3\}\}}
   (which corresponds to no verifiable record) by the verified entry
   \texttt{@article\{BurnsFlach2001Doc,
   author = \{Burns, D. and Flach, M.\},
   title = \{Tamagawa numbers for motives with (non-commutative) coefficients\},
   journal = \{Documenta Mathematica\}, volume = \{6\}, pages = \{501--570\},
   year = \{2001\}, url = \{https://eudml.org/doc/123971\}\}}.
\item Add the entry
   \texttt{@book\{EGM1998, author = \{Elstrodt, J. and Grunewald, F. and Mennicke, J.\},
   title = \{Groups Acting on Hyperbolic Space: Harmonic Analysis and Number Theory\},
   series = \{Springer Monographs in Mathematics\}, publisher = \{Springer\},
   year = \{1998\}, isbn = \{9783540627456\}\}}.
\item Add the entry
   \texttt{@article\{Lichnerowicz1963, author = \{Lichnerowicz, A.\},
   title = \{Spineurs harmoniques\},
   journal = \{C. R. Acad. Sci. Paris\}, volume = \{257\}, pages = \{7--9\},
   year = \{1963\}\}}.
\item Add the entry
   \texttt{@article\{Castella2024, author = \{Castella, F.\},
   title = \{Tamagawa number conjecture for CM modular forms\},
   note = \{arXiv:2407.11891\}, year = \{2024\}\}}.
\item Add the entry
   \texttt{@article\{FlachSurvey2004, author = \{Flach, M.\},
   title = \{The Equivariant Tamagawa Number Conjecture: A survey\},
   note = \{Expanded version of Baltimore 2002 talk; with appendix by C. Greither\},
   year = \{2004\}, url = \{https://www.pma.caltech.edu/documents/5558/baltimore-final.pdf\}\}}.
\item Add the entry
   \texttt{@incollection\{Deligne1979, author = \{Deligne, P.\},
   title = \{Valeurs de fonctions L et p\'eriodes d'int\'egrales\},
   booktitle = \{Automorphic forms, representations and L-functions (Corvallis 1977), Part 2\},
   series = \{Proc. Sympos. Pure Math.\}, volume = \{XXXIII\},
   publisher = \{Amer. Math. Soc.\}, year = \{1979\}, pages = \{313--346\}\}}.
\end{enumerate}

\noindent All cited arXiv IDs (\texttt{1503.04626}, \texttt{2407.11891})
are pre-verified in \cite[\S M.G.2]{BIGTABLE2026}.  The Burns--Flach
correction is anti-fab catch \#45 (internal, pre-publication).

% End of sec_07_5_A3_Lichnerowicz.tex
